\section{The Deployment Front-Door Rule}
\label{sec:rule}

The attribution gap calls for a limited information rule rather than a new general law of AI liability. The rule must tell an affected person whom to approach, tell enterprises what record to maintain, and tell courts when and how that record becomes available. It must also provide a prompt exit for an enterprise whose connection to the brand or technical stack did not include practical control over the asserted risk.

This Part supplies that architecture. Section A ties coverage to sectoral legislative designation. Section B makes the visible operator the first-answer gateway and defines a threshold linked to an existing legal claim. Section C gives \term{material control} a cumulative, risk-specific test. Section D separates record creation, preservation, and staged production. Section E specifies confidentiality, remedies, and safe exit. Section F consolidates the proposal into model statutory text, and Section G explains how an ex ante state-law record duty can operate alongside the Federal Rules.

\subsection{Sectoral Designation Defines the Field}
\label{sec:designation}

``High impact'' should identify a deployment category, not attach permanently to a model. A general-purpose model can support casual drafting in one setting and become part of a benefits, hiring, clinical, or physical-control process in another. The deployment supplies the protected interest, affected population, operating institution, and legally relevant control. Coverage should follow those features.

The proposed statute therefore applies only when a legislature has designated a category and specified six matters: the protected interests; the visible operator; the minimum deployment record; the retention period and material-change trigger; the claimant's threshold; and the available enforcement remedy. A legislature can make the designation directly or authorize a sectoral agency to maintain a bounded list under statutory criteria. The legislature retains the choices that determine who must answer and which interests justify compulsory recordkeeping.

Five features identify the strongest candidates. First, the deployment affects an interest already protected by substantive law, such as employment opportunity, access to public benefits, health, credit, housing, education, physical safety, or the welfare of minors. Second, the system repeatedly mediates decisions or exposure rather than producing an isolated low-consequence convenience. Third, affected people predictably observe less about the operative system than the operator and suppliers know. Fourth, the identity, version, control, and event fields can be standardized at reasonable cost. Fifth, existing doctrine can adjudicate the dispute once those facts are available. These criteria select the settings in which a record gateway unlocks an established legal regime rather than inventing a new one.

Designation also fixes the public-law interface. If a covered operator is a government institution, the legislature should specify the enforcing tribunal, available relief, applicable exhaustion route, and the treatment of sovereign or official immunity. The front door need not imply a damages action against the government. It can supply an administrative record, declaratory remedy, reviewable denial, or other access path tailored to the underlying interest. Defining that route at designation keeps the information right aligned with the sector's remedial structure.

This design avoids an unstable universal definition. Employment screening, allocation of essential public benefits, clinical diagnosis or treatment, credit, housing, education admissions, physical safety, and relational services offered to minors may justify different records and retention periods. A companion service may require age-access, relational-behavior, crisis-intervention, and continuity records. An employment tool may require input-feature, validation, selection-rate, configuration, and version records. A clinical system may require intended-use, validation-population, uncertainty, escalation, override, and follow-up records. The shared front-door architecture permits sector-specific content.

California's enacted companion-chatbot provisions demonstrate the feasibility of this form. The legislature defined a covered relational service and its operator, required specified disclosures and self-harm protocols, imposed protections for known minors and reporting obligations, supplied a private remedy, and preserved cumulative law.\lrfootnote{See \statcite{californiaSB2432025}. The enacted provisions do not contain the deployment front-door rule proposed here; they demonstrate that a legislature can designate a relational deployment, identify its operator, and attach enforceable protocol and reporting duties.} The statute did not need to declare every conversational model high impact. It regulated a deployment category and the enterprise that made it available.

A designation should also identify exclusions. Research or internal evaluation that does not produce a legally or practically consequential effect on an affected person ordinarily falls outside. So does general cloud hosting without risk-specific configuration or intervention authority. Small entities may receive simplified record forms or regulator assistance where the sector permits, but the information essential to identify a consequential decision should not disappear with enterprise size. A covered operator using open weights still controls a deployment; an upstream author who has no continuing authority may fall outside the post-release control lane.

\subsection{The Visible Operator Answers First}
\label{sec:gateway}

The rule resolves the service problem by selecting the entity an affected person can ordinarily identify. The \term{visible operator} is the enterprise or public institution that made the covered deployment available to the affected person or incorporated it into the decision process that affected that person. It may be a companion-service provider, employer, hospital, lender, benefits agency, school, or operator of a branded application. Its gateway role follows from visibility and provenance, not a presumption that it controlled every risk.

The operator receives a verified front-door notice. The notice must identify: (1) a legally cognizable injury or an existing entitlement to equitable relief; (2) the cause of action, statutory right, or protected legal interest the claimant expects to invoke; (3) concrete facts supporting a plausible nexus between the asserted risk and a designated deployment; (4) the relevant time, transaction, account, decision, or interaction with as much specificity as reasonably available; and (5) the limited identity and record categories sought. ``An AI harmed me'' is insufficient. The facts must connect the deployment to a legal interest and connect the requested record to an element, defense, or remedial question. Verification means a sworn declaration or an unsworn declaration under penalty of perjury based on personal knowledge, with specifically identified matters permitted on information reasonably available to the requester. Counsel separately certifies the request under the forum's ordinary good-faith standard.

The notice invokes a statutory access right that exists outside merits discovery. It can be served before a complaint, with a complaint, or after filing. Pre-suit service gives the operator an opportunity to answer without motion practice and preserves the claimant's ability to name the correct parties. Service should follow a designated electronic or registered-agent channel. The operator receives thirty days, subject to one short extension for documented technical or privacy needs, to answer, object by record category, or state that a named deployment did not process the identified event. Silence, a blanket proprietary designation, and referral to an unnamed vendor are not responses.

If the operator refuses or materially underanswers, the claimant may bring a narrow enforcement action against that known entity. The pleaded statutory violation is the qualifying deployment, compliant request, service, and nonperformance; the claimant need not plead the still-concealed merits facts to enforce the access right. The administering agency or a court with an independent basis for jurisdiction can order the defined first answer. The claimant bears the burden to establish coverage and a qualifying notice. The operator bears the burden to support an objection based on nonoperation, burden, privilege, privacy, security, or the absence of a requested field. Jurisdiction, immunity, preemption, and a pure legal bar that can be resolved on the request remain available before any access beyond identity and preservation.

The gateway must provide a \term{first answer}. It identifies the application and model versions active for the relevant event; the upstream and downstream entities that supplied, configured, integrated, or operated material components; the entity that holds each listed record category; and the persons or roles with authority to update, restrict, roll back, suspend, or stop the relevant function. It preserves records within its control, identifies the retention schedule and any known deletion, and gives statutory notice to each enterprise that may have occupied a material control lane. It also provides the incident-specific output, decision, notice, configuration, warning, and intervention events within its custody, subject to privilege and lawful protection of third-party data.

This first answer solves two practical problems. The claimant receives the deployment's identity and provenance without reconstructing the operator's vendor architecture. Potential control holders receive timely notice and can preserve their own records. A visible operator cannot answer that the vendor relationship is confidential or that its procurement team does not know which version ran; maintaining that provenance is part of the ex ante deployment-record duty.

The visible operator remains free to establish that its gateway role did not include material control over the asserted risk. A hospital that only offered physical access to an independently governed research system, a marketplace that processed payment without configuring the application, or an employer whose vendor acted outside any retained or delegated hiring function may present that record. Visibility opens the door; the test in the next Section determines who remains.

\subsection{Material Control Requires Risk, Authority, and Consequence}
\label{sec:material-control}

``Ability to affect the system'' is too broad. A cloud provider, investor, data vendor, researcher, model author, and downstream institution may each have some causal connection to a deployment. The front door should reach an enterprise only when three requirements are satisfied at the relevant time.

First, \term{risk specificity}: the asserted control must connect to the risk alleged in the verified notice. General influence over corporate strategy, compute availability, or model research does not establish control over age access, hiring rank, clinical escalation, a particular tool permission, or another specified hazard.

Second, \term{practical authority}: the enterprise must have possessed an operational or legally enforceable ability to configure, constrain, evaluate, monitor, update, suspend, roll back, or terminate the relevant function. A theoretical technical capacity, an expired contract right, or influence that required another actor's voluntary cooperation is evidence but not practical authority by itself. Contractual audit, update, safety-override, and suspension rights matter because they may be usable controls; a contractual label declaring the parties independent does not decide the question.

Third, \term{materiality}: exercising the authority must have been capable of making a substantial difference to the risk's creation, exposure, detection, or reduction. An incidental contribution or remote ability affecting all customers in the same undifferentiated way falls outside. Materiality does not require proof that the intervention would certainly have prevented the claimant's injury. It requires a nonmarginal control lane relevant to the deployment and risk.

The requirements are cumulative. Nonexclusive factors include who selected the deployment's purpose and population; chose the model and version; configured instructions, memory, retrieval, data fields, tools, and credentials; set engagement, ranking, decision, or escalation thresholds; received evaluations, complaints, and incident signals; could issue an update or impose a restriction; retained contractual audit or safety-override rights; and could suspend, roll back, or terminate the function. The inquiry is temporal. Authority obtained after the event does not create earlier control, although later records may reveal what authority existed.

Several negative rules keep the test bounded. Investment, ordinary cloud capacity, publication of research, past employment, trademark association, and model provenance alone do not establish material control. Publishing open weights may end the original provider's ability to monitor or stop an independent deployment. A downstream modifier that removes safeguards, adds credentials, or selects a new high-impact purpose can occupy the relevant lane instead. Conversely, a provider that retains mandatory updates, remote suspension, risk monitoring, or binding deployment conditions may retain a lane despite contract language assigning ``all responsibility'' downstream.

Material control determines information obligations, not substantive liability. An enterprise can possess a control relevant to feasibility yet owe no duty to the claimant under governing law. It can receive an alert that falls short of legally sufficient notice. The test identifies the entities and records needed to ask the merits questions accurately. The statute's \term{merits firewall} preserves that separation: front-door participation and record production create no inference on an element of the underlying claim.

\subsubsection{Adjudicating Material Control}

The test becomes useful doctrine when each requirement has an evidentiary question and an allocated burden. After the first answer, the claimant identifies a concrete lane that the record connects to the asserted risk. The claimant might point to an employer's enabled ranking feature, a provider's retained update authority, an application's self-harm threshold, or an integrator's control of a workflow trigger. The showing is about capability and allocation, not ultimate breach. An enterprise seeking to terminate its front-door role then identifies the record that defeats risk specificity, practical authority, or materiality. The tribunal decides the issue on the closed Stage One record or permits a bounded Stage Two inquiry when a genuine conflict remains.

\begin{table}[htbp]
\centering
\caption{Material Control as an Adjudicable Test}
\label{tab:material-control-proof}
\small
\begin{tabularx}{\textwidth}{>{\raggedright\arraybackslash}p{0.18\textwidth} >{\raggedright\arraybackslash}p{0.35\textwidth} X}
\toprule
Requirement & Evidence that supports the lane & Evidence that supports termination of the lane \\
\midrule
Risk specificity & Function map; risk ownership; feature requirements; evaluation scope; incident category; warnings; escalation assignment. & General corporate influence; unrelated research; infrastructure service without risk-specific configuration; records addressing a different feature or hazard. \\
Practical authority & Contract right; administrative permission; release or update role; audit access; monitoring feed; restriction, rollback, suspension, or termination mechanism; consistent exercise in practice. & Expired or legally unavailable right; access requiring another actor's uncommitted consent; technical ability prohibited and unused; independent modification beyond visibility or reach. \\
Materiality & Control-path record showing the action could alter exposure, detection, decision weight, access, response, or duration; tests or operations showing a nonmarginal effect. & Feature inactive for the event; output never generated or transmitted; control unrelated to the asserted path; action incapable of changing the identified risk. \\
\bottomrule
\end{tabularx}
\end{table}

Reserved contractual rights deserve particular care. A right to obtain validation on demand, require a safety update, audit use, or suspend a feature is practical authority when the right was legally available and operationally usable. The holder need not have exercised it in the claimant's event; otherwise an enterprise could erase its information obligation by leaving a reserved safeguard idle. A generic right to terminate an entire commercial agreement may be too remote when it provides no timely, risk-specific intervention. The contract supplies evidence, while administrative permissions, actual practice, response time, technical dependencies, and past interventions show whether the right could function.

Actual power can matter when a contract is silent or inaccurate. A vendor that routinely pushes mandatory updates, disables a feature across tenants, monitors risk signals, or supplies the only version map may possess practical authority even if the agreement calls the customer solely responsible. A customer that can enable, ignore, or combine a score may control the decision lane even when the vendor hosts the data. A systems integrator with one-time access during installation may fall outside after handoff unless it retained monitoring, credentials, or change authority. The inquiry reconstructs the operative deployment rather than enforcing labels selected for another purpose.

Materiality imposes the counterfactual discipline. The question is whether exercise of the identified authority could substantially alter creation, exposure, detection, or reduction of the asserted risk. It is not whether that exercise would certainly have prevented the completed injury. Stage One can answer many cases: the feature was disabled; no score was generated; the output never reached the decisionmaker; the provider had no post-release connection; or the supposed controller joined the stack after the event. Where the event path and authority conflict, Stage Two can examine configuration history, tests, monitoring, or a representative protocol limited to that lane.

The same test produces different allocations across architectures. In an integrated companion service, the visible operator may control access, relationship design, memory, engagement objectives, warnings, monitoring, and account restriction; an infrastructure provider offering generic compute exits. In an employment stack, the employer may control job criteria, feature activation, decision weight, review, and final disposition, while a vendor controls model design, defaults, validation, updates, and cross-customer monitoring. In an open-weights deployment, the original publisher may have no continuing control after release while a modifier and operating institution acquire the risk-specific tools, data, credentials, and stop authority. The rule follows each lane without requiring every participant to remain.

A material-control ruling terminates only the special front-door role. An entity never named on an underlying claim returns to ordinary nonparty status after completing any first-answer duty. A party already facing an independently sufficient cause of action remains subject to that claim even if it lacks the control lane alleged in the notice. This distinction gives safe exit its intended economy without using an information proceeding to decide a merits question it was not designed to reach.

\subsection{Creation, Preservation, and Production Are Separate Duties}
\label{sec:three-duties}

The front door depends on three legal functions that begin at different times.

\subsubsection{Record Creation Before Operation}

The \term{deployment-record duty} arises before a covered deployment affects the public. The operator must create and maintain a record identifying the system, version, purpose, population, material components, vendors, data sources, tools, credentials, intended limits, risk owners, release approval, and intervention authority. During operation, it must record material changes, evaluations, complaints, incidents, escalations, overrides, restrictions, rollbacks, and termination decisions. An enterprise with a retained control lane must maintain the subset corresponding to that lane and provide the operator with current provenance and safety-contact information.

A divided deployment requires a retrieval architecture as well as a list of vendors. Before use, the visible operator should obtain contractual rights to retrieve the identity, version, configuration, event, and control-allocation fields assigned to an upstream holder on demand. The agreement should identify a responsive agent, stable deployment and event identifiers, notice and preservation duties, lawful routes for producing customer-controlled data, successor obligations after acquisition, and the treatment of subcontractors. The operator need not warehouse a supplier's source code or unrelated customer data. It must keep a closed local record sufficient to identify what ran and a legally usable path to obtain the defined supplier fields.

That requirement converts discovery doctrine into deployment design. Where Rule 34 control turns on a legal right to obtain documents on demand, a covered operator can create that right before deployment instead of litigating it after harm. Where privacy law or a customer agreement places legal control with the operating institution, the contract can authorize a coordinated response rather than allowing the host and customer to direct the requester to each other. A foreign supplier remains subject to ordinary jurisdiction and service limits, but a domestic operator can preserve the minimum record locally and condition covered use on a reachable agent and retrieval duty.

The designation supplies the detail. A statute need not require indefinite retention of every prompt or every user's data. It should define event categories, sampling where appropriate, privacy minimization, and a retention period calibrated to the protected interest and ordinary limitations period. High-impact decisions affecting employment or benefits may require the inputs, output, version, governing policy, notice, and review event for each affected person while permitting aggregation or deletion of unrelated interaction content. Relational deployments involving minors may require crisis detections and interventions while minimizing intimate content unrelated to the asserted event.

\subsubsection{Preservation After a Defined Trigger}

Retention and litigation preservation are related but distinct. The statute fixes a baseline period before any dispute. A verified front-door notice, regulator demand, filed claim, or other event making a dispute reasonably foreseeable suspends ordinary deletion for the risk-specific record. The operator must issue a hold to relevant custodians and notify identified control holders. The notice should not freeze every model and customer record. Its scope follows the event, versions, period, user or decision, control lanes, and asserted risk.

This statutory trigger creates a stable bridge to existing preservation law. It also protects defendants: a dated notice and documented hold define what had to be preserved and when. The record can show that a missing item was outside scope, deleted before any duty attached, or available from another source.

\subsubsection{Staged Production After the Threshold}

Production proceeds in stages. Stage One is the first answer described above: identity, provenance, active versions, relevant entities and custodians, allocation of operational and stop authority, incident-specific events, policies and warnings in force, preservation status, and control provisions in contracts. Its purpose is to identify the deployment and test the control lanes.

Stage Two follows only when Stage One and the verified notice support a claim-specific need. It can include relevant evaluations and validation; complaints and incident trends; configuration history; release and change approvals; rejected or conditional mitigations; escalation and response; and records needed to test a proposed alternative or defense. Time, version, affected population, feature, and risk category bound the request.

Source code, model weights, unrestricted training data, security-sensitive controls, and unrelated users' content occupy a final tier. A court should order them only when the requesting party shows that a material issue cannot reasonably be tested through less intrusive evidence such as versioned documentation, validation results, representative testing, interrogatories, an expert protocol, or in camera review. The distinction protects legitimate secrecy without converting a trade secret into a barrier to the legal process that governs its consequential use.

\subsection{Protection, Remedies, and Safe Exit}
\label{sec:remedies}

Confidentiality is part of the rule's design. Production remains limited to nonprivileged matter. Courts may use redaction, pseudonymization, data minimization, attorneys'-eyes-only access, secure expert environments, in camera review, and protective orders for trade secrets, research, security, medical information, and personal data.\lrfootnote{See \statcite{frcp26}, r. 26(c)(1)(G).} The operator must still identify that a protected record exists, its custodian, date, version, and asserted basis for withholding. A privilege log separates a lawful protection from an unreviewable claim that the record is proprietary.

The statute should specify graduated consequences.

\begin{enumerate}[leftmargin=2.2em]
\item \term{Late or incomplete identification}: an order enforcing the statutory first answer and reasonable expenses caused by noncompliance. For a claim created or governed by the same enactment, the limitations period is suspended from service of a compliant request until a fixed period after compliance or final denial. Other state and federal claims retain their own limitations and tolling law.
\item \term{Failure to create or retain the minimum record}: the administrative penalty, corrective record order, compliance monitoring, or statutory-access remedy specified by the sectoral enactment. The consequence addresses the independent deployment obligation; it does not become a prescribed inference on the underlying injury.
\item \term{Loss after a preservation trigger}: an order to identify the loss and available substitute sources, followed by measures under the governing forum's preservation law. For electronically stored information that should have been preserved in anticipation or conduct of federal litigation and cannot be restored or replaced, Rule 37(e) supplies the federal standard. The state statute should not prescribe a competing litigation sanction for the same loss.
\item \term{Intentional concealment or falsification}: specified civil penalties, enforcement expenses, and remedies for the independent access violation. Litigation misconduct remains subject to the tribunal's governing authority.
\item \term{The merits firewall}: compliance, noncompliance, and front-door participation create no presumption of duty, breach, defect, discriminatory operation, agency, causation, damages, or joint liability. Governing law determines whether a record omission has any separate relevance to an independently pleaded claim.
\end{enumerate}

Safe exit is the reciprocal limit. After Stage One, an enterprise may move to terminate its front-door role by showing that it lacked material control over the asserted risk during the relevant period. The claimant may use the first-answer record to contest that showing. An entity not otherwise sued returns to ordinary nonparty status. A party remains only when an independently sufficient claim, a valid subpoena, or regulatory authority supplies a basis for continued participation. The gateway obligation therefore ends at the point where the control record shows that another lane or no AI-specific lane is relevant.

\subsection{Model Statutory Text}
\label{sec:model-text}

The following text consolidates the rule. Sectoral legislation would replace the bracketed specifications and conform remedies to the governing field.

\begin{quote}
\small
\textbf{Section 1. Designated Deployment.} This Act applies only to a category of AI deployment designated by [the legislature] and identified by its covered function, protected interests, visible operator, affected persons, minimum deployment record, retention period, material-change threshold, and administering authority.

\textbf{Section 2. Deployment Record.} Before making a covered deployment available or incorporating it into a decision process, a visible operator shall create and maintain a record identifying the system and versions; purpose and affected population; material model, data, retrieval, tool, and credential components; suppliers and integrators; risk and release owners; evaluation and approval; material limitations and unresolved risk decisions; incident and escalation channels; and authority to update, restrict, roll back, suspend, or terminate. Each enterprise retaining material control shall maintain the corresponding portion and provide current provenance and safety-contact information to the operator. The record shall be updated upon a material change and retained for [period].

\textbf{Section 3. Verified Front-Door Notice.} An affected person may serve the visible operator's designated agent with a sworn declaration or declaration under penalty of perjury stating a legally cognizable injury or existing entitlement to equitable relief; the legal claim or protected interest to which the requested record relates; specific facts supporting a plausible and risk-specific nexus to the covered deployment; the relevant event, transaction, account, application, decision, or period with as much specificity as reasonably available; and the closed identity and record fields sought. Counsel, if any, shall certify that the request has a proper purpose and factual basis after reasonable inquiry.

\textbf{Section 4. First Answer and Preservation.} Within [thirty] days after service, the operator shall provide the application and model versions active for the event; each entity that supplied, configured, integrated, or operated a material component; the custodian and legal controller of each requested record category; the event-specific output or decision and whether it was transmitted or used; and each entity with update, restriction, rollback, suspension, or termination authority. The operator shall preserve its risk-specific record, identify known deletion or unavailability, provide the remaining Stage One fields designated for the sector, and notify each entity that may possess material control. A notified entity shall preserve its corresponding record and acknowledge the lane it holds or the basis on which it disclaims a lane. One extension of [fifteen] days may be obtained for stated technical, privacy, or security need.

\textbf{Section 5. Material Control.} An enterprise possesses material control only when, at the relevant time, its authority was (a) specifically connected to the asserted risk; (b) practically exercisable through configuration, constraint, evaluation, monitoring, update, restriction, rollback, suspension, termination, or a legally enforceable equivalent; and (c) capable of making a substantial difference to the risk's creation, exposure, detection, or reduction. Investment, general cloud service, publication, research, brand association, provenance, or theoretical technical capacity alone is insufficient.

\textbf{Section 6. Objection and Enforcement.} The operator may object by specific record field and shall state the facts supporting noncoverage, nonoperation, burden, privilege, privacy, security, or unavailability. A blanket proprietary designation or referral to an unnamed supplier does not satisfy this Section. An affected person may bring an action to enforce Sections 3 and 4 in [the designated tribunal]. The requester bears the burden to establish coverage, service, and a qualifying notice. The operator bears the burden on a field-specific objection based on facts within its control. The tribunal may allow limited-purpose participation by a notified enterprise and shall resolve the request on an expedited record.

\textbf{Section 7. Staged Access and Protection.} The minimum first answer is an access obligation arising from covered deployment. Further nonprivileged production may be obtained by agreement or ordered in a pending proceeding only when the notice and initial record establish claim-specific need, and shall be limited by time, version, affected population, feature, and asserted risk. Source code, weights, sensitive data, and security controls require a showing that less intrusive evidence cannot reasonably resolve a material issue. The enforcing tribunal shall protect privilege, privacy, trade secrets, research, and security through redaction, minimization, pseudonymization, limited access, secure expert review, in camera review, or another lawful safeguard.

\textbf{Section 8. Merits Firewall and Termination.} Compliance, noncompliance, notice, and participation create no presumption of duty, breach, defect, discriminatory operation, agency, causation, damages, or joint liability. Governing law supplies every element and defense. An enterprise that completes Section 4 and establishes that it lacked material control may terminate its front-door role. An entity not otherwise subject to an independently sufficient claim resumes nonparty status. A valid subpoena, regulatory authority, or another source of law may require continued participation on its own terms.

\textbf{Section 9. Remedies and Limitations.} The tribunal may order compliance and award reasonable expenses caused by nonperformance. Failure to create or retain the minimum record is subject to [specified administrative or civil remedies] for the independent deployment violation. Intentional concealment or falsification is subject to [specified civil penalties]. For a claim created or governed by this Act, the applicable limitations period is suspended from service of a compliant request until [sixty] days after compliance or final denial. This Section does not alter the limitations period or tolling rule governing an independent federal claim or a claim beyond the enacting jurisdiction's authority.

\textbf{Section 10. Nonwaiver and Procedural Coordination.} The duties to create, retain, and provide the minimum deployment record are nonwaivable as to affected persons, arise from covered deployment, and exist independently of anticipated or pending litigation. Nothing in this Act alters jurisdiction, immunity, privilege, or the Federal Rules of Civil Procedure governing pleading, service, discovery, amendment, protective orders, or sanctions. In a federal action, any additional production and any measure for lost electronically stored information proceed under otherwise applicable federal procedure. Rights and remedies supplied by other law remain available on their own terms.
\end{quote}

The text makes the product legible. The regulated acts are operating a covered deployment without the minimum record and refusing the resulting access right, not unsafe operation in the abstract. The trigger is tied to a legal claim and concrete deployment facts. Control is risk specific. Access is staged. Confidential information receives protection. A record holder can terminate its special role. The merits remain where ordinary law placed them.

\subsection{Federal Procedure Implements Rather Than Displaces the Duty}
\label{sec:erie}

A state front-door statute should regulate deployment conduct rather than attempt to rewrite federal litigation. The minimum-record duty attaches when the operator makes a designated deployment available or incorporates it into a covered decision. The minimum-access duty attaches when an affected person serves a qualifying request. Both exist whether a lawsuit is anticipated, pending, arbitrated, or never filed. The operator's statutory violation is complete when it fails to create, retain, or provide the defined record on the enactment's terms.

That structure separates the right from its enforcement forum. A stand-alone access action requires an independent basis for subject-matter jurisdiction; the state statute does not manufacture federal jurisdiction. A claimant may proceed in the designated state forum, invoke diversity when its requirements are met, assert a federal front-door right if Congress enacts one, or seek supplemental jurisdiction when the access claim forms part of the same case or controversy as a claim properly in federal court. Once the dispute is in federal court, federal rules govern the manner of pleading, service, discovery, amendment, protection, and sanctions.

The Supreme Court's framework makes direct collision the first question. \emph{Hanna} applies a valid Federal Rule when it answers the question before the court. \emph{Burlington Northern} likewise refused a state appellate penalty that occupied the same field as a federal rule governing the consequence at issue. \emph{Walker}, by contrast, applied a state service-linked limitations rule because Rule 3 did not purport to answer when the state limitations period was tolled.\lrfootnote{See \casecite[464--65, 468, 471--74]{hanna1965}; \casecite[4--8]{burlingtonNorthern1987}; \casecite[749--52]{walker1980}; \statcite{rulesEnablingAct}.} The front door is designed for the \emph{Walker} side of that line: it defines what a covered operator must have recorded and what minimum information an affected person may demand before litigation. It neither advances the Rule 26 discovery timetable nor enlarges Rule 34 control or Rule 45 subpoena power.

The allocation then follows by function. State law can require a version, provenance, control, and event record years before a dispute; define a nonwaivable request; authorize an action for refusal; and suspend the limitations period for claims the state enactment governs. Rule 8 controls a complaint filed in federal court. Rules 26, 34, and 45 control additional discovery and its sequence. Rule 26(c) supplies federal protective orders. Rule 15 controls amendment and relation back where it applies. The state statute gives the parties an existing object and an access entitlement; the Federal Rules provide the litigation machinery.

Preservation requires the same separation. The deployment statute defines a baseline retention period and can penalize noncreation or premature deletion as an operating violation. A qualifying request can also create an independent statutory hold. Rule 37(e), however, governs federal litigation measures when electronically stored information that should have been preserved in anticipation or conduct of the action is lost and cannot be restored or replaced. The Rule supplies its own prejudice and intent standards; it does not create the underlying preservation duty.\lrfootnote{See \statcite{frcp37e}.} The model text therefore specifies administrative or access remedies for the deployment violation and leaves federal adverse-inference and case-management measures to Rule 37(e) within its scope.

Limitations follows the legislature's authority. A state enactment can suspend the period for a state claim it creates or governs while the visible operator supplies the identity record. It should not promise tolling for an independent federal cause of action or automatic relation back against a later-named party. Congress can pair a federal front-door right with federal tolling where national uniformity is warranted. In either form, the access period gives the claimant time to use the answer while preserving Rule 15's role in a later federal case.

The result is a clean institutional division. Deployment law determines that a consequential system must remain identifiable and grants a defined access right to the person it affects. Federal procedure governs how a federal court adjudicates enforcement and any later merits claim. The front door arrives with the missing record; it does not arrive as a substitute code of civil procedure.
